% Created 2026-03-05 Thu 00:16
% Intended LaTeX compiler: pdflatex
\documentclass[11pt]{article}
\usepackage[utf8]{inputenc}
\usepackage[T1]{fontenc}
\usepackage{graphicx}
\usepackage{longtable}
\usepackage{wrapfig}
\usepackage{rotating}
\usepackage[normalem]{ulem}
\usepackage{amsmath}
\usepackage{amssymb}
\usepackage{capt-of}
\usepackage{hyperref}
\author{The Prologos Project - from Zachary Larson}
\date{2026-03-05}
\title{Propagator Networks\\\medskip
\large A Comprehensive Survey of Theory, Practice, and Applications}
\hypersetup{
 pdfauthor={The Prologos Project - from Zachary Larson},
 pdftitle={Propagator Networks},
 pdfkeywords={},
 pdfsubject={},
 pdfcreator={Emacs 30.2 (Org mode 9.7.39)}, 
 pdflang={English}}
\begin{document}

\maketitle
\tableofcontents

\section{Introduction}
\label{sec:org1511226}

A \textbf{propagator network} is a model of computation in which autonomous,
stateless computational agents --- called \textbf{propagators} --- communicate
through shared stateful storage elements --- called \textbf{cells} --- that
accumulate partial information about values. Unlike traditional models
where programs execute as sequential instructions over mutable memory,
propagator networks describe computation as the gradual, multidirectional
accumulation of knowledge, driven by local rules, until no further
deductions can be made.

The key insight, due to Alexey Radul \footnote{Radul, A. (2009). "Propagation Networks: A Flexible
and Expressive Substrate for Computation." PhD thesis, MIT.\label{org978162e}}, is that \textbf{a cell
should not be seen as storing a value, but as accumulating information
about a value}. This reframing transforms cells from simple memory
locations into lattice-structured accumulators of partial information,
and elevates the computational model from "compute a result" to "refine
what we know until we know enough."

Propagator networks unify ideas from constraint propagation, truth
maintenance, dataflow programming, and lattice-theoretic fixed-point
computation into a single, composable framework. They have applications
in type inference, constraint solving, abstract interpretation,
incremental compilation, reactive programming, and protocol
verification. Despite this breadth, the core model is remarkably simple:
cells hold lattice elements; propagators are monotone functions between
them; and the network runs until quiescence (a fixed point).

There is, as of this writing, no Wikipedia article on propagator
networks. This document aims to fill that gap: a self-contained survey
accessible to a reader with general programming background, progressing
from foundations through theory to applications and the state of the
art.
\section{Origins and History}
\label{sec:orgfcaf9d2}

\subsection{Constraint Propagation in Circuit Analysis (1977)}
\label{sec:org6eeffdb}

The propagator model traces its lineage to Gerald Jay Sussman and
Richard Stallman's 1977 work on computer-aided circuit analysis
\subsection{Steele's Constraint Language (1980)}
\label{sec:org9b00a1c}

Guy Lewis Steele Jr., also working under Sussman at MIT, developed the
theoretical foundations of constraint-based programming in his 1980 PhD
thesis \footnote{Steele, G.L. (1980). "The Definition and
Implementation of a Computer Programming Language Based on
Constraints." PhD thesis, MIT. AI Laboratory TR-595.}. His key contribution was the \textbf{locality
principle}: each constraint operates using only locally available
information and places newly derived values on locally attached wires.
No constraint needs global knowledge of the network; all computation is
local.

This principle became a cornerstone of the propagator model. It is what
makes propagator networks inherently parallelizable: if every
computational element operates independently using only local
information, there is no need for global coordination.
\subsection{Truth Maintenance Systems (1979--1986)}
\label{sec:orgd9c48ee}

Sussman's circuit analysis work spawned a line of research into \textbf{truth
maintenance} --- systems that track \emph{why} facts are believed, not just
\emph{what} is believed.

\textbf{Jon Doyle's TMS (1979)} \footnote{Doyle, J. (1979). "A Truth Maintenance System."
\emph{Artificial Intelligence}, 12:231--272.} introduced the
justification-based truth maintenance system. Each belief in the system
is annotated with its justification --- the set of prior beliefs and
inference rules that support it. When a supporting belief is retracted,
all conclusions that depended on it are automatically invalidated. This
enabled non-monotonic reasoning: the system can change its mind as
assumptions change.

\textbf{Johan de Kleer's ATMS (1986)} \footnote{de Kleer, J. (1986). "An Assumption-Based TMS."
\emph{Artificial Intelligence}, 28(2):127--162.} advanced the idea
significantly with the assumption-based truth maintenance system.
Rather than maintaining one "current" set of beliefs (as the TMS does),
the ATMS simultaneously tracks \emph{all possible} assumption sets --- each a
consistent "worldview." With \(n\) assumptions, there are potentially
\(2^n\) environments. Since all environments are maintained in parallel,
switching between worldviews is free (no retraction or re-derivation
needed).

The ATMS proved essential for the propagator model's handling of
hypothetical reasoning and search, as we will see.
\subsection{Waltz Filtering and Arc Consistency (1975--1977)}
\label{sec:org0d36c57}

Independently from the Sussman lineage, David Waltz's 1975 algorithm for
interpreting line drawings of 3D scenes \footnote{Waltz, D. (1975). "Understanding Line Drawings of Scenes
with Shadows." In \emph{The Psychology of Computer Vision}, McGraw-Hill.} demonstrated that
local constraint propagation could dramatically reduce search spaces.
Alan Mackworth formalized this as \textbf{arc consistency} in 1977
\subsection{Dataflow Programming (1974)}
\label{sec:orgdd3763f}

Jack Dennis formalized dataflow schemas in 1974 \footnote{Dennis, J.B. (1974). "First Version of a Data Flow
Procedure Language." \emph{Lecture Notes in Computer Science}, vol 19,
Springer.}: programs
as directed graphs where operations execute upon data availability,
rather than following a sequential program counter. Gilles Kahn
independently introduced Kahn Process Networks \footnote{Kahn, G. (1974). "The Semantics of a Simple Language for
Parallel Programming." \emph{IFIP Congress}, pp. 471--475.} --- sequential
processes communicating through FIFO queues --- and proved that such
networks are deterministic: the output sequence does not depend on
process execution speed.

Propagator networks generalize dataflow in two ways: (1) information
flows \emph{multidirectionally} (not just producer-to-consumer), and (2)
cells accumulate \emph{partial information} (not complete values). A standard
dataflow network is a special case of a propagator network where all
propagators are unidirectional and all cells hold complete values.
\subsection{Radul's PhD Thesis (2009)}
\label{sec:orgff987da}

Alexey Radul's 2009 PhD thesis at MIT, "Propagation Networks: A Flexible
and Expressive Substrate for Computation" \textsuperscript{\ref{org978162e}}, supervised
by Sussman, is the definitive formulation of the propagator model as a
general-purpose computational paradigm. The thesis synthesized three
decades of ideas --- constraint propagation, truth maintenance, lattice
theory, dataflow --- into a unified model with a precise mathematical
foundation.

Radul argued that conventional programming systems are insufficiently
expressive because they are built on the image of a single computer with
global effects on a large memory. The propagator model offers an
alternative: autonomous machines communicating through shared cells,
where computation is the progressive accumulation of partial information
about values.

The thesis demonstrated that an astonishing range of computational
paradigms embed naturally into propagator networks: expression
evaluation, constraint satisfaction, truth maintenance, functional
reactive programming, rule-based systems, and type inference are all
instances of the same underlying pattern.
\subsection{"The Art of the Propagator" (2009)}
\label{sec:org524c90e}

Radul and Sussman published a companion paper, "The Art of the
Propagator" \footnote{Radul, A. \& Sussman, G.J. (2009). "The Art of the
Propagator." MIT-CSAIL-TR-2009-002.}, presenting the model more concisely. The
paper emphasizes that the model "makes it easy to smoothly combine
expression-oriented and constraint-based programming, and easily
accommodates implicit incremental distributed search in ordinary
programs."
\subsection{The Revised Report (2010)}
\label{sec:org427cf9c}

Radul and Sussman published the "Revised Report on the Propagator Model"
\subsection{Software Design for Flexibility (2021)}
\label{sec:orgd35cf7c}

Chris Hanson and Sussman's book \emph{Software Design for Flexibility}
\section{Core Concepts}
\label{sec:org7d03620}

\subsection{Cells}
\label{sec:org56cd361}

A \textbf{cell} is the stateful storage element of a propagator network. Unlike
a variable in traditional programming (which holds a single definitive
value and can be overwritten), a cell \emph{accumulates information about a
value}. A cell supports three operations:

\begin{enumerate}
\item \textbf{Add content}: New information is \emph{merged} with the cell's existing
information (not overwritten).
\item \textbf{Read content}: Retrieve the current state of accumulated knowledge.
\item \textbf{Register propagators}: Propagators subscribe to the cell and are
notified when its content changes.
\end{enumerate}

The information in a cell ranges from "I know absolutely nothing" (the
special value \texttt{nothing}, or lattice bottom \(\bot\)) to "I know everything
there is to know, and the answer is 42" to the anomalous case "I know
there is a contradiction!" (lattice top \(\top\)).

Critically, information in a cell \emph{never decreases}. A cell can move
from "I don't know" to "it's an integer" to "it's a positive integer" to
"it's 42," but it can never move backward from "it's 42" to "I don't
know." This monotonicity is not merely a convention --- it is the
mathematical foundation ensuring convergence and determinism.
\subsection{Propagators}
\label{sec:orgc1045d1}

A \textbf{propagator} is the computational element of the network. Propagators are:

\begin{itemize}
\item \textbf{Autonomous}: They operate independently, with no central controller
directing their execution.
\item \textbf{Stateless}: They maintain no internal state. All persistent state
resides in cells.
\item \textbf{Asynchronous}: They can fire in any order; the final result is the
same (confluence).
\end{itemize}

A propagator monitors its input cells and, when any input changes,
computes new information and adds it to its output cells. The
computation is always \emph{additive} --- a propagator can only \emph{contribute}
information, never \emph{remove} it.

In practice, a propagator is a function from the current network state to
a (possibly modified) network state: it reads from some cells, computes,
and writes to others via the merge operation. If its contribution adds
no new information (i.e., the merge result is identical to the cell's
existing content), the network state is unchanged and no further
propagation is triggered.
\subsection{The Merge Operation}
\label{sec:orgc826b69}

The merge operation is the heart of the propagator model. When a
propagator writes information \(b\) to a cell already containing
information \(a\), the cell's new content is \(\text{merge}(a, b)\) --- the
combination of both pieces of knowledge.

Merge must satisfy three algebraic laws:

\begin{itemize}
\item \textbf{Commutativity}: \(\text{merge}(a, b) = \text{merge}(b, a)\)
--- the order in which information arrives does not matter.
\item \textbf{Associativity}: \(\text{merge}(\text{merge}(a, b), c) = \text{merge}(a, \text{merge}(b, c))\)
--- grouping does not matter.
\item \textbf{Idempotency}: \(\text{merge}(a, a) = a\)
--- re-contributing the same information has no effect.
\end{itemize}

These three laws make merge a \textbf{join-semilattice} operation (see
\ref{sec:org063aa5a}). The commutativity and associativity
ensure that the final cell state is independent of the order in which
propagators fire. The idempotency ensures that redundant contributions
are harmless. Together, these properties guarantee \textbf{confluence}: the
result of running a propagator network depends only on the initial cell
contents and the set of propagators, not on the order or timing of
execution.

Key behaviors of merge:

\begin{itemize}
\item \(\text{merge}(\bot, x) = x\) --- merging anything with "nothing" yields
that thing (bottom is the identity).
\item \(\text{merge}(x, x) = x\) --- consistent redundant information is
harmless.
\item \(\text{merge}(\top, x) = \top\) --- once a contradiction is reached, it
cannot be undone (top is absorbing).
\item \(\text{merge}(a, b) = \top\) when \(a\) and \(b\) are incompatible ---
contradictory information produces a contradiction.
\end{itemize}
\subsection{Partial Information}
\label{sec:org05317c2}

The propagator model supports a rich hierarchy of partial information types:

\begin{itemize}
\item \textbf{Nothing (\(\bot\))}: Complete ignorance about the value.
\item \textbf{Raw values}: Fully determined information (e.g., the integer 42).
\item \textbf{Intervals}: Bounded ranges (e.g., \([3.0, 5.0]\)), representing the
knowledge that a value lies somewhere within bounds.
\item \textbf{Finite domains}: Sets of possible values (e.g., \(\{1, 2, 3, 4\}\)).
\item \textbf{Supported values}: Values paired with the set of assumptions that
justify them (see \ref{sec:org6c85cbb}).
\item \textbf{TMS entries}: Sets of supported values, representing multiple
contingent beliefs that coexist.
\end{itemize}

Each partial information type defines three generic operations:
\texttt{equivalent?} (test whether two structures represent the same
information), \texttt{merge} (combine two structures), and \texttt{contradictory?}
(test whether a structure represents an impossible state).

The power of the propagator model comes from the ability to define \emph{new}
partial information types. Any data type that forms a join-semilattice
with a bottom element and a contradiction test can be used as cell
content. This extensibility is what allows propagator networks to
operate over types, multiplicities, session protocols, intervals, finite
domains, or any other structured form of partial knowledge.
\subsection{Monotonicity}
\label{sec:orgd4073be}

All information flow in a propagator network must be \textbf{monotonic}: cells
can only gain information, never lose it. In lattice-theoretic terms,
the content of a cell can only move \emph{upward} in the information ordering
(toward more-determined states, further from \(\bot\), closer to \(\top\)).

Monotonicity is the property that makes everything else work:

\begin{itemize}
\item \textbf{Confluence}: Because merge is commutative and associative, and
propagators can only increase cell values, the final result is
independent of execution order.
\item \textbf{Termination}: In a finite lattice, information can only increase a
finite number of times before reaching a fixed point. (Infinite
lattices require \emph{widening} --- see \ref{sec:orgddf2bf4}.)
\item \textbf{Compositionality}: Networks can be freely combined without worrying
about interference. Each sub-network's monotonic contributions compose
with others.
\end{itemize}
\subsection{Quiescence and Fixed Points}
\label{sec:org9157938}

A propagator network reaches \textbf{quiescence} when no propagator has new
information to contribute --- the system has reached a \textbf{fixed point}.
At quiescence, every cell contains the \textbf{least fixed point} of the system
of equations defined by the propagators: the smallest (least-informative)
cell contents consistent with all propagators being satisfied.

The existence and uniqueness of this least fixed point is guaranteed by
the \textbf{Knaster-Tarski fixed-point theorem}: every monotone function on a
complete lattice has a least fixed point, and it can be computed by
iterating from \(\bot\):

$$\text{lfp}(f) = \bigsqcup_{n \in \mathbb{N}} f^n(\bot)$$

This is precisely what running a propagator network to quiescence
computes: starting from cells containing \(\bot\) (nothing), repeatedly
firing propagators until nothing changes.
\subsection{Multidirectional Computation}
\label{sec:org2691539}

Unlike traditional functions (which map inputs to outputs in one
direction), propagators support \textbf{multidirectional computation}.
A constraint like \(x + y = z\) is expressed not as a function from
\((x, y)\) to \(z\), but as a set of propagators:

\begin{itemize}
\item One computes \(z\) from \(x\) and \(y\)
\item One computes \(x\) from \(z\) and \(y\)
\item One computes \(y\) from \(z\) and \(x\)
\end{itemize}

Given any two of the three values, the third is determined. Given all
three, consistency is verified. This arises naturally from the merge
semantics: each propagator contributes whatever it can, and merge
reconciles all contributions.

The classic example (from Radul and Sussman) is temperature conversion.
The relationship \(C = (F - 32) \times 5/9\) is expressed once, and the
network automatically solves for whichever variable is unknown ---
forward (F to C) or backward (C to F).
\subsection{Supported Values and Dependencies}
\label{sec:org6c85cbb}

A \textbf{supported value} is a pair \((v, P)\) of a datum \(v\) and a set of
\textbf{premises} \(P\) --- the assumptions on which \(v\) depends. This enables
the system to track not just \emph{what} it believes, but \emph{why}.

When a propagator derives new information, the result is tagged with the
union of its inputs' premise sets. If cell \(A\) holds \((3, \{p_1\})\) and
cell \(B\) holds \((5, \{p_2\})\), an addition propagator writes \((8,
\{p_1, p_2\})\) to cell \(C\) --- "the sum is 8, given that we believe both
\(p_1\) and \(p_2\)."

When a contradiction is discovered, the system can identify \emph{which
premises} contributed. The conflicting premise set is recorded as a
\textbf{nogood} --- a set of assumptions that cannot all be simultaneously true.
This enables \textbf{dependency-directed backtracking}: the system retracts the
minimal set of assumptions needed to resolve the conflict, rather than
blindly reverting the most recent choice.
\subsection{Truth Maintenance Integration}
\label{sec:org9c4b2c5}

Propagator networks integrate with \textbf{truth maintenance systems} to enable
hypothetical reasoning. A TMS-backed cell contains a set of supported
values, each contingent on different assumptions. When assumptions are
activated or deactivated, the TMS automatically updates which values are
"in" (currently believed) and which are "out."

The ATMS (Assumption-Based Truth Maintenance System) extends this to
maintain \emph{all possible} worldviews simultaneously. Rather than committing
to one set of assumptions, the ATMS tracks every consistent combination.
This enables:

\begin{itemize}
\item \textbf{Hypothetical reasoning}: "What if we assume \(p_1\)? What follows?"
--- without committing to \(p_1\).
\item \textbf{Dependency-directed backtracking}: When a contradiction is found,
trace back to the responsible premises and record the nogood.
\item \textbf{Non-monotonic reasoning}: While the underlying information flow is
monotonic (lattice joins), the TMS layer enables \emph{belief revision} by
changing which premises are active.
\end{itemize}
\section{Lattice-Theoretic Foundations}
\label{sec:org063aa5a}

\subsection{Lattices and Partial Orders}
\label{sec:org288f4b7}

A \textbf{partial order} is a set \(S\) with a binary relation \(\leq\) that is
reflexive (\(a \leq a\)), antisymmetric (\(a \leq b\) and \(b \leq a\)
implies \(a = b\)), and transitive (\(a \leq b\) and \(b \leq c\) implies \(a
\leq c\)). The ordering represents "has at least as much information as."

A \textbf{join-semilattice} \((L, \leq, \bot, \sqcup)\) is a partial order with:

\begin{itemize}
\item A \textbf{bottom element} \(\bot\) (least informative --- "nothing")
\item A \textbf{join operation} \(\sqcup\) (least upper bound --- "combine
information")
\end{itemize}

The join is commutative, associative, and idempotent. Join is precisely
the merge operation of propagator cells.

A \textbf{complete lattice} additionally has a \textbf{top element} \(\top\) (maximally
informative, representing contradiction) and a \textbf{meet} operation \(\sqcap\)
(greatest lower bound). In propagator networks, reaching \(\top\) signals
that incompatible information has been merged.
\subsection{Merge as Lattice Join}
\label{sec:org3046f89}

The merge operation implements the lattice join:

$$\text{merge}(a, b) = a \sqcup b$$

All algebraic properties of merge follow from the lattice axioms:

\begin{center}
\begin{tabular}{lll}
Property & Equation & Consequence\\
\hline
Commutativity & \(a \sqcup b = b \sqcup a\) & Order-independence\\
Associativity & \((a \sqcup b) \sqcup c = a \sqcup (b \sqcup c)\) & Grouping-independence\\
Idempotency & \(a \sqcup a = a\) & Redundancy-tolerance\\
Identity & \(a \sqcup \bot = a\) & Adding nothing changes nothing\\
Annihilation & \(a \sqcup \top = \top\) & Contradiction is irreversible\\
\end{tabular}
\end{center}
\subsection{Monotone Functions and Convergence}
\label{sec:orgf7abde4}

A function \(f : L \to L\) is \textbf{monotone} if \(a \leq b\) implies \(f(a) \leq
f(b)\) --- more information in, at least as much information out.
Propagators must be monotone functions.

The \textbf{Knaster-Tarski fixed-point theorem} guarantees that every monotone
function on a complete lattice has a least fixed point. Combined with
Kleene's iteration theorem, this means:

\begin{enumerate}
\item A fixed point exists.
\item It can be computed by iterated application from \(\bot\).
\item For finite lattices, iteration terminates in finitely many steps.
\item The result is independent of the iteration order.
\end{enumerate}

Property (4) --- \emph{confluence} --- is the deepest guarantee of the
propagator model: no matter how propagators are scheduled, the network
converges to the same result.
\subsection{The CALM Theorem}
\label{sec:org526044c}

The CALM theorem (Consistency As Logical Monotonicity) \footnote{Hellerstein, J.M. \& Alvaro, P. (2019). "Keeping CALM:
When Distributed Consistency Is Easy." \emph{Communications of the ACM},
63(9).\label{orge067b9c}},
formulated by Hellerstein and Alvaro, establishes that monotonic
programs (those expressible as monotone functions on lattices) can be
executed in a distributed, coordination-free manner while still
producing consistent results. Non-monotonic operations (like negation or
aggregation) require coordination (synchronization barriers).

Propagator networks are inherently monotonic, so the CALM theorem
guarantees that they can be distributed across multiple processors or
machines without coordination, and the result will be the same as
single-machine execution. This makes propagator networks a natural fit
for parallel and distributed computation.
\subsection{Widening and Narrowing}
\label{sec:orgddf2bf4}

For lattices of infinite height (e.g., the lattice of integer intervals
\([lo, hi]\)), Kleene iteration may not terminate --- a cell could ascend
through infinitely many refinements. \textbf{Widening} (due to Cousot \& Cousot
\section{Relationship to Other Computational Paradigms}
\label{sec:org28167ef}

\subsection{Constraint Programming}
\label{sec:orgea00c65}

Propagator networks subsume the constraint propagation phase of
constraint programming (CP) systems. In CP, variables have domains
(typically finite sets), and constraints are propagators that narrow
domains by removing inconsistent values. Modern production CP solvers
like Gecode \footnote{Schulte, C. \& Stuckey, P.J. (2008). "Efficient Constraint
Propagation Engines." \emph{ACM Transactions on Programming Languages and
Systems}, 31(1).\label{org4786136}} are explicitly built on propagator architectures.

The relationship is direct:

\begin{center}
\begin{tabular}{ll}
Constraint Programming & Propagator Networks\\
\hline
Variable domain & Cell content (lattice value)\\
Constraint propagator & Propagator (fire function)\\
Arc consistency (AC-3, etc.) & Run to quiescence\\
Backtracking search & ATMS worldview branching\\
Learned nogood & Recorded assumption conflict\\
\end{tabular}
\end{center}

Propagator networks generalize CP in that cells can hold \emph{any}
lattice-structured partial information, not just finite domains. An
interval, a type, a session protocol, or a user-defined lattice are all
valid cell contents.
\subsection{Datalog and Fixed-Point Logic}
\label{sec:orgfcd42f3}

Datalog --- the database query language based on Horn clauses evaluated
to a fixed point --- is a direct instance of propagator-style
computation on a specific lattice: sets of tuples ordered by subset
inclusion.

\begin{itemize}
\item Each Datalog relation is a cell holding a set of tuples.
\item Each Datalog rule is a propagator: when input relations change,
derive new tuples for the head relation.
\item Naive evaluation (iterate all rules until nothing changes) =
run to quiescence.
\item Semi-naive evaluation (only propagate \emph{new} tuples) = an optimization
of the propagator scheduling strategy.
\end{itemize}

Flix \footnote{Madsen, M., Yee, M.-H., \& Lhotak, O. (2016). "From Datalog
to Flix: A Declarative Language for Fixed Points on Lattices." \emph{PLDI 2016}.}, a research language by Madsen et al. (PLDI 2016),
extends Datalog by replacing set inclusion with arbitrary user-defined
lattices. Flix's model is essentially Datalog-as-propagator-network:
lattice-valued relations, monotone rules, iterative fixed-point
evaluation. Datafun \footnote{Arntzenius, M. \& Krishnaswami, N.R. (2016). "Datafun: A
Functional Datalog." \emph{ICFP 2016}.} (Arntzenius \& Krishnaswami, ICFP 2016)
takes this further, embedding monotone fixed-point computation into a
functional type system.
\subsection{Reactive Programming and Spreadsheets}
\label{sec:orgd9706a2}

Spreadsheets are the most widely deployed example of propagator-like
computation. When a cell changes, dependent cells recompute
automatically --- unidirectional propagation through a dependency graph.

Reactive programming generalizes this with observable streams and
automatic change propagation. Functional reactive programming (FRP)
adds compositional abstractions.

Propagator networks differ from reactive/FRP systems in three key ways:

\begin{enumerate}
\item \textbf{Multidirectionality}: Reactive systems propagate change in one
direction (from sources to dependents). Propagators support
multidirectional constraints.
\item \textbf{Partial information}: FRP cells hold complete values that are
overwritten on change. Propagator cells accumulate partial
information via merge.
\item \textbf{Merge semantics}: FRP cells are "last writer wins." Propagator cells
are "all writers contribute" via lattice join.
\end{enumerate}

Radul's thesis demonstrates that FRP embeds naturally within the
propagator framework: a standard reactive system is a propagator network
with unidirectional propagators and complete-value cells. David
Thompson's work \footnote{Thompson, D. (2014). "Functional Reactive User
Interfaces with Propagators."} shows the practical value of this
embedding for building functional reactive user interfaces with
propagators.
\subsection{SAT/SMT Solvers}
\label{sec:org67eee2a}

The relationship between propagators and SAT/SMT solvers is structural.
Modern SAT solvers use CDCL (Conflict-Driven Clause Learning):

\begin{enumerate}
\item \textbf{Decision}: Choose an unassigned variable and assign it.
\item \textbf{Unit propagation}: Deduce forced assignments from clauses.
\item \textbf{Conflict analysis}: When a contradiction is found, learn a new
clause (nogood).
\item \textbf{Non-chronological backjumping}: Jump back to the most relevant
decision level.
\end{enumerate}

Each of these maps directly to propagator concepts:

\begin{center}
\begin{tabular}{ll}
CDCL & Propagator Network\\
\hline
Decision literal & ATMS assumption\\
Unit propagation & Run to quiescence\\
Conflict clause & Nogood set\\
Non-chronological backjump & Dependency-directed backtracking\\
Theory solver & Domain-specific propagators\\
\end{tabular}
\end{center}

SMT (Satisfiability Modulo Theories) extends SAT with "theory
propagators" that participate in propagation and conflict analysis ---
the same architectural pattern as domain-specific propagators in a
propagator network.

D'Silva et al. \footnote{D'Silva, V. et al. (2013). "Abstract Conflict Driven
Learning." \emph{SAS 2013}.\label{org393d4e3}} showed that CDCL can be lifted from the
Boolean lattice to arbitrary lattice-based abstract domains ("Abstract
CDCL"). This directly connects SAT/SMT solving to the lattice-theoretic
core of propagator networks.
\subsection{The Actor Model}
\label{sec:orgbac87a9}

Carl Hewitt's Actor Model \footnote{Hewitt, C. (1973). "A Universal Modular ACTOR Formalism
for Artificial Intelligence."} describes computation as
autonomous agents (actors) communicating through asynchronous messages.
Propagator networks resemble actor systems architecturally --- both
feature autonomous computational elements without central control ---
but differ crucially:

\begin{center}
\begin{tabular}{lll}
Dimension & Actors & Propagators\\
\hline
Communication & Point-to-point messages & Shared cells with merge\\
State & Per-actor mutable state & Shared cells accumulate info\\
Determinism & Non-deterministic & Deterministic (confluence)\\
Convergence & No guarantee & Guaranteed (lattice theory)\\
Fault model & Supervision, let-it-crash & Contradiction + ATMS nogood\\
\end{tabular}
\end{center}

The lattice constraint on cell values is what buys determinism: unlike
actors, where message ordering can affect the outcome, the commutativity
of merge ensures that propagator execution order is irrelevant.

Kuper and Newton's LVars \footnote{Kuper, L. \& Newton, R.R. (2013). "LVars: Lattice-based Data
Structures for Deterministic Parallelism." \emph{FHPC 2013}.} (see ) formalize this connection: LVars are shared mutable
variables with monotonic writes and threshold reads, guaranteeing
deterministic parallelism --- precisely the semantics of propagator
cells.
\subsection{CRDTs and Distributed Systems}
\label{sec:org561d394}

\textbf{Conflict-Free Replicated Data Types} (CRDTs) \footnote{Shapiro, M. et al. (2011). "Conflict-Free Replicated Data
Types." \emph{SSS 2011}.} share the
same mathematical substrate as propagator cells. State-based CRDTs
(CvRDTs) require:

\begin{itemize}
\item A join-semilattice structure on state
\item A merge function that is commutative, associative, and idempotent
\item Monotonic state updates
\end{itemize}

These are \emph{exactly} the requirements for propagator cell contents.
Both CRDTs and propagators solve the problem of combining information
from multiple autonomous sources in a consistent, order-independent
way. CRDTs solve it for distributed replicas across a network;
propagators solve it for computational elements within a program.

The CALM theorem \textsuperscript{\ref{orge067b9c}} makes this connection precise: monotonic
computations (like propagator networks and CRDT merges) can be
distributed without coordination. Edward Kmett's influential talks on
propagators \footnote{Kmett, E. (2016--2019). "Propagators." Talks at FnConf,
YOW!, and others.} explicitly identify this unification: "propagators,
CRDTs, Datalog, SAT solving, and FRP all share the same
lattice-fixpoint structure."
\section{Applications}
\label{sec:orga9bab48}

\subsection{Type Inference}
\label{sec:orgc741aaf}

Type inference --- the problem of determining the types of expressions
in a program without explicit annotations --- maps directly onto
propagator networks. Radul's thesis includes a chapter titled "Type
Inference Looks Like Propagation Too."

The mapping:

\begin{itemize}
\item Each type variable (metavariable) is a \emph{cell}, initially containing
\(\bot\) ("I don't yet know what type this is").
\item Each typing rule (function application, variable binding, etc.) is a
\emph{propagator} that constrains related type cells.
\item Unification is \emph{merge}: when two type cells are constrained to be
equal, their contents are merged (unified). Compatible types merge to
a more-specific type; incompatible types merge to \(\top\)
(type error).
\item Running to quiescence computes the principal type (the most general
type consistent with all constraints).
\end{itemize}

Bidirectional type checking --- splitting inference into synthesis
(bottom-up, from term to type) and checking (top-down, from expected
type to term) --- is a natural instance of multidirectional propagation.
Recent work by Leijen \footnote{Leijen, D. (2025). "Omnidirectional Type Inference
for ML."} on "omnidirectional type
inference" makes this explicit: typing constraints suspend when
information is insufficient and resume when other constraints supply it
--- precisely propagator semantics.
\subsection{Constraint Solving}
\label{sec:org41de211}

Propagator networks are the native architecture of constraint solving.
Variables are cells, constraints are propagators, and solving is running
to quiescence. The advantages over traditional constraint solvers:

\begin{itemize}
\item \textbf{Mixed-domain solving}: Finite domain constraints, interval
constraints, type constraints, and user-defined constraints can
coexist in a single network, connected by cross-domain propagators.
\item \textbf{Persistent backtracking}: If the network uses persistent (immutable)
data structures, backtracking is \(O(1)\) --- just keep the old network
reference. Traditional CP systems pay \(O(\text{trail length})\) per
backtrack.
\item \textbf{Extensibility}: Adding a new constraint domain is defining a new
lattice and wiring propagators. In traditional systems (Gecode,
MiniZinc), adding a domain requires low-level implementation.
\end{itemize}
\subsection{Abstract Interpretation}
\label{sec:org9053461}

Abstract interpretation \footnote{Cousot, P. \& Cousot, R. (1977). "Abstract
Interpretation: A Unified Lattice Model for Static Analysis of Programs
by Construction or Approximation of Fixpoints." \emph{POPL 1977}.} --- computing sound
approximations of program behavior over abstract domains --- is
propagator networks operating over abstract lattices:

\begin{itemize}
\item Each \emph{program point} is a cell holding an abstract value (e.g., sign,
interval, pointer alias set).
\item Each \emph{program statement} is a propagator computing abstract transfer
functions.
\item Running to quiescence computes the analysis fixed point.
\item Widening handles infinite-height abstract domains.
\end{itemize}

The connection to propagators is made precise by \textbf{Galois connections}:
pairs of monotone functions \((\alpha : C \to A, \gamma : A \to C)\)
connecting concrete and abstract domains. Cross-domain propagators use
Galois connections to flow information between multiple abstract domains
simultaneously --- the "reduced product" of multiple analyses becomes
automatic when all domains share a propagator network.
\subsection{Protocol Verification via Session Types}
\label{sec:orgabb0b61}

Session types --- type-theoretic specifications of communication
protocols --- can be verified using propagator networks. Each channel
endpoint is a cell in a session type lattice; each communication
operation (send, receive, select, offer) is a propagator that constrains
the channel's protocol.

Caires and Pfenning \footnote{Caires, L. \& Pfenning, F. (2010). "Session Types
as Intuitionistic Linear Propositions." \emph{CONCUR 2010}.} established a Curry-Howard
correspondence between session types and intuitionistic linear logic.
Propagator-based session type checking exploits this: the duality
constraint (client and server must follow dual protocols) is a
bidirectional propagator, and protocol violations are contradictions
(lattice top) with dependency chains explaining \emph{which} communication
step violated \emph{which} protocol clause.
\subsection{Incremental Compilation}
\label{sec:orgc79f174}

An incremental compiler can be structured as a propagator network:

\begin{itemize}
\item Source files are cells (content = parsed AST).
\item Compilation phases (parse → elaborate → type-check → codegen) are
propagator chains.
\item A file change updates a cell; running to quiescence recomputes only
affected downstream cells.
\item Persistent networks give \(O(1)\) access to the previous version for
delta computation.
\end{itemize}

The multidirectionality of propagators adds a novel capability: a type
signature change can propagate \emph{backward} to callers, enabling "what-if"
analysis ("what happens if I change this type?") without recompiling.
This connects to Acar's work on self-adjusting computation
\subsection{Diagnostic Reasoning}
\label{sec:orgddc88f8}

De Kleer and Williams's General Diagnostic Engine (GDE, 1987)
\section{Implementations}
\label{sec:orgf1dc383}

\subsection{The Reference Implementation (Scheme)}
\label{sec:org7769237}

The canonical implementation is Radul and Sussman's Scheme-Propagators
\subsection{Haskell Libraries}
\label{sec:org9ae21d5}

\textbf{Edward Kmett's \texttt{propagators}} \footnote{Kmett, E. Haskell \texttt{propagators} library.
\url{https://github.com/ekmett/propagators}}: An exploration
of propagator design in Haskell. The primary innovation beyond the
published Scheme work is the use of \textbf{observable sharing} to let users
write in a direct programming style that is automatically transformed
to and from propagator form.

\textbf{Holmes} \footnote{Holmes (i-am-tom). Haskell constraint-solving with
propagators and CDCL. \url{https://github.com/i-am-tom/holmes}}: A reference library for constraint-solving with
propagators and CDCL. Holmes exposes two strategies: \texttt{satisfying}
(returns first valid configuration) and \texttt{whenever} (returns all valid
configurations). Its README demonstrates a complete Sudoku solver.
\subsection{Clojure Libraries}
\label{sec:org62e9c38}

\textbf{propaganda} \footnote{Nilsson, T.K. \texttt{propaganda} --- Clojure propagator
library. \url{https://github.com/tgk/propaganda}}: A Clojure propagator library implementing
the model from "The Art of the Propagator." Offers both an STM-based
and an immutable implementation strategy.
\subsection{Constraint Programming Systems}
\label{sec:orgcef0018}

\textbf{Gecode} \textsuperscript{\ref{org4786136}}: A production-quality, highly efficient constraint
solver explicitly built on a propagator architecture. Gecode
demonstrates that the propagator architecture is viable for
industrial-strength constraint solving. Its propagation kernel is
domain-independent, and propagators are first-class objects that can be
dynamically added to a constraint model.

Most modern CP systems (SICStus Prolog, ECLiPSe, SWI-Prolog's CLP
libraries) use propagator-based architectures internally, where each
constraint is a propagator that narrows variable domains.
\section{A Worked Example: The Prologos Propagator Architecture}
\label{prologos}
To illustrate how propagator networks function as practical
infrastructure, we describe the propagator architecture of Prologos ---
a programming language that uses propagator networks as the unified
substrate for type inference, multiplicity tracking (QTT), and session
type verification.
\subsection{Design Principles}
\label{sec:org08bf018}

Prologos's propagator network is built on three architectural principles:

\begin{enumerate}
\item \textbf{Persistent/immutable data structures}: The entire network is a pure,
immutable value. Each operation (cell creation, propagator addition,
cell write) returns a \emph{new} network, leaving the old one unchanged.
This enables \(O(1)\) backtracking (keep the old reference) and safe
parallel exploration (fork the network value).

\item \textbf{Three-domain unification}: Type inference, QTT multiplicity
tracking, and session type verification all operate on the same
propagator network. Each domain defines its own lattice and
propagators; cross-domain \emph{Galois connection} propagators bridge them.

\item \textbf{Trait-based extensibility}: Lattice behavior (bottom, join,
contradiction detection) is defined via a \texttt{Lattice} trait.
New domains can be added at the library level by implementing
the trait, without modifying the propagator infrastructure.
\end{enumerate}
\subsection{Core Data Structures}
\label{sec:org91b37b5}

The network is built on CHAMP (Compressed Hash Array Mapped Prefix-tree)
persistent hash maps, giving \(O(\log_{32} N)\) lookup, insert, and delete
with structural sharing.

\begin{verbatim}
struct cell-id(n : Nat)          ;; unique cell identifier
struct prop-id(n : Nat)          ;; unique propagator identifier

struct prop-cell
  value      : Lattice           ;; current lattice element (starts at ⊥)
  dependents : CHAMP(prop-id)    ;; propagators to fire on change

struct propagator
  inputs  : List(cell-id)        ;; cells this reads from
  outputs : List(cell-id)        ;; cells this may write to
  fire-fn : Network → Network    ;; pure state transformer

struct prop-network
  cells         : CHAMP(cell-id → prop-cell)
  propagators   : CHAMP(prop-id → propagator)
  worklist      : List(prop-id)  ;; propagators waiting to fire
  fuel          : Nat            ;; step limit to prevent runaway
  contradiction : cell-id | #f   ;; first cell that reached ⊤
  merge-fns     : CHAMP(cell-id → (a b → merged))
  ...                            ;; + contradiction-fns, widen-fns,
                                 ;;   decomposition registries
\end{verbatim}
\subsection{The Three Lattice Domains}
\label{sec:org5375384}

\subsubsection{Type Lattice}
\label{sec:orgbcdf05f}

\begin{verbatim}
type-bot (⊥, no information — fresh metavariable)
    ↓
T (concrete type expression)
    ↓
type-top (⊤, contradiction — incompatible types)
\end{verbatim}

The merge function attempts \emph{pure} unification (side-effect free) of
two type expressions. If unification succeeds, the result is the unified
type. If it fails, the result is \(\top\) (type error).

A critical design choice: the type lattice's merge function uses only
\emph{read-only callbacks} to resolve metavariables, breaking the circular
dependency between the metavar store (which uses the network) and the
type lattice (which the network uses).
\subsubsection{Multiplicity Lattice}
\label{sec:org2136a55}

\begin{verbatim}
mult-bot (⊥, no information)
    ↓
m0 / m1 / mw (erased, linear, unrestricted — mutually incomparable)
    ↓
mult-top (⊤, contradiction)
\end{verbatim}

This is a flat lattice: \(\bot \sqcup x = x\), \(x \sqcup x = x\), \(x
\sqcup y = \top\) for \(x \neq y\). Multiplicities track how many times a
value is used, supporting Quantitative Type Theory (QTT).
\subsubsection{Session Lattice}
\label{sec:org3afce1c}

\begin{verbatim}
sess-bot (⊥, no information)
    ↓
Send(A, S) / Recv(A, S) / Choice{...} / Offer{...} / End / ...
    ↓
sess-top (⊤, protocol violation)
\end{verbatim}

Session types describe communication protocols. Merge performs structural
unification: same-polarity sessions (send + send) merge component types
and continuations; different-polarity sessions (send + recv) produce
\(\top\) (protocol violation). Choice labels merge covariantly (intersection);
offer labels merge contravariantly (union).
\subsection{Cross-Domain Propagation via Galois Connections}
\label{sec:org31cc2e6}

Different lattice domains are connected by \textbf{Galois connection
propagators} --- pairs of monotone functions \((\alpha, \gamma)\) that
bridge information between domains:

\begin{verbatim}
;; Type → Multiplicity bridge:
;;   α extracts multiplicity info from type structure
;;   γ returns ⊥ (types don't learn from multiplicities)

;; Session → Type bridge:
;;   α extracts message type from session step
;;   γ returns ⊥ (sessions don't learn from types)
\end{verbatim}

Each Galois connection creates two unidirectional propagators (one per
direction). The \(\alpha\) direction (concrete → abstract) extracts
information; the \(\gamma\) direction typically returns \(\bot\) (the
abstract domain doesn't constrain the concrete domain). More
sophisticated connections can propagate information in both directions.
\subsection{Scheduling Strategies}
\label{sec:org2e3aa88}

\subsubsection{Gauss-Seidel (Sequential)}
\label{sec:org4df115a}

The default scheduler processes one propagator at a time:

\begin{verbatim}
while worklist ≠ empty:
  pop propagator p from worklist
  apply fire-fn(p) to network, producing network'
  if any output cell changed: enqueue dependent propagators
  decrement fuel
\end{verbatim}

This is the simplest and typically fastest strategy for single-threaded
execution.
\subsubsection{BSP (Parallel)}
\label{sec:org0876f56}

The Bulk Synchronous Parallel scheduler processes all worklist entries in
parallel against a frozen snapshot, then bulk-merges results:

\begin{verbatim}
while worklist ≠ empty:
  Round k:
    fire ALL propagators in worklist (in parallel, against snapshot)
    bulk-merge all cell changes
    enqueue affected propagators for next round
\end{verbatim}

By the CALM theorem, both strategies converge to the same fixed point.
The BSP strategy can leverage multiple cores but may require more rounds
(Jacobi iteration converges more slowly than Gauss-Seidel for chain
topologies).
\subsection{Structural Decomposition}
\label{sec:org0d88e7b}

When the network unifies two compound types (e.g., \(\text{Pi}(A_1,
B_1)\) with \(\text{Pi}(A_2, B_2)\)), the unification propagator creates
\emph{sub-cells} for the components (\(A_1\) vs \(A_2\), \(B_1\) vs \(B_2\)) and
registers new sub-propagators to unify them. A deduplication registry
prevents the same compound pair from being decomposed twice.

This on-demand structural decomposition is key to handling recursive and
compound types efficiently: cells and propagators are created only as
needed, keeping the network sparse.
\subsection{Example: Type Inference as Propagation}
\label{sec:org1dea322}

Consider type-checking the expression \texttt{f x} where \texttt{f : Int -> Bool} and
\texttt{x : Int}:

\begin{enumerate}
\item Create cells: \(c_f\) (type of \texttt{f}), \(c_x\) (type of \texttt{x}), \(c_r\)
(type of result).
\item Write known types: \(c_f \leftarrow \text{Int} \to \text{Bool}\), \(c_x
   \leftarrow \text{Int}\).
\item Add application propagator: "if \(c_f = A \to B\) and \(c_x\) is
compatible with \(A\), then \(c_r \leftarrow B\)."
\item Propagator fires: \(c_f\) is \(\text{Int} \to \text{Bool}\), \(c_x\) is
\(\text{Int}\). Merge \(c_x\) with domain type \(A = \text{Int}\):
consistent. Write \(c_r \leftarrow \text{Bool}\).
\item Quiescence: all cells stable. Result type is \(\text{Bool}\).
\end{enumerate}

If \texttt{x} had been \texttt{"hello"} (type \(\text{String}\)), step 4 would merge
\(\text{String}\) with \(\text{Int}\), producing \(\top\) (type error) ---
with a dependency chain explaining that the error arose from the
application of \texttt{f} to \texttt{x}.
\subsection{Propagators as Universal Substrate}
\label{sec:org1cb94ce}

The deepest validation of this architecture is that \emph{the same
infrastructure serves all three domains without modification}. The
propagator network does not know or care whether a cell holds a type, a
multiplicity, or a session protocol. It only knows that cells have merge
functions, propagators have fire functions, and the network runs to
quiescence.

This means that the three domains compose for free: type inference,
multiplicity tracking, and session type verification run simultaneously
on the same network, with Galois connections bridging information
between them. A type error can surface from a session type constraint;
a multiplicity violation can arise from a type structure. The
propagator network routes information to wherever it is needed,
automatically.
\section{Future Directions}
\label{sec:org56f8ccf}

\subsection{Theorem Proving via Abstract CDCL}
\label{sec:org822417e}

D'Silva et al. \textsuperscript{\ref{org393d4e3}} showed that CDCL generalizes from Booleans to
arbitrary lattices. Combined with an ATMS, this yields a
propagator-native theorem prover: cells hold lattice values, propagators
encode inference rules, ATMS manages hypothetical reasoning and conflict
analysis. This could enable type-level theorem proving and automated
capability verification without an external SMT solver.
\subsection{User-Defined Constraint Domains}
\label{sec:org512d634}

Because lattice behavior is defined by a trait, users of a
propagator-based language can define their own constraint domains
entirely at the library level. A string-length-bound lattice, a
nullability lattice, or a permission lattice can be defined, populated
with propagators, and connected to the type system via Galois
connections --- all without modifying the core infrastructure.
\subsection{Distributed Propagation}
\label{sec:org0557a49}

Since propagator cells satisfy CRDT merge requirements, a propagator
network can be partitioned across nodes. Boundary cells are replicated
and synchronized via lattice merge. Message ordering is irrelevant
(commutativity), and convergence is guaranteed (CALM theorem). This
enables distributed type checking, multi-agent constraint solving, and
decentralized belief propagation.
\subsection{Incremental and Live Programming}
\label{sec:orgbb8c7c4}

Persistent propagator networks enable "what-if" exploration: fork the
network, make a hypothetical change, observe consequences, and discard
the fork at zero cost. Combined with IDE integration, this enables live
type checking where edits propagate incrementally to diagnostics,
completions, and refactoring suggestions.
\subsection{Provenance and Explainability}
\label{sec:org8e2e4d0}

ATMS dependency chains provide causal explanations for every derived
value. Combined with provenance semirings \footnote{Green, T.J., Karvounarakis, G., \& Tannen, V.
(2007). "Provenance Semirings." \emph{PODS 2007}.},
this enables rich answers to questions like "why does this cell have
this type?" "what assumptions contributed?" and "what is the minimal
change to resolve this error?" --- transforming compiler diagnostics
from symptoms to explanations.
\section{References}
\label{sec:org91a08b8}
\end{document}
